% 保留原结构图的数据流、训练关系及执行控制关系，仅调整节点与走线。
% policystage 决定当前方法包含的模块，policyfocus 决定派生图的强调部分。
\begin{tikzpicture}[
  font=\small,
  >={Stealth[length=2mm]},
  box/.style={rounded corners=2pt,align=center,inner sep=2pt,outer sep=0pt},
  data/.style={box,draw=black!60,fill=white,minimum width=28mm,text width=26mm,minimum height=10mm},
  process/.style={box,draw=blue!60!black,fill=blue!6},
  retained/.style={box,draw=blue!55!black,fill=blue!4},
  offline/.style={box,draw=gray!65,fill=gray!7},
  added/.style={draw=orange!85!black,fill=orange!10,line width=1pt},
  container/.style={draw=blue!50!black,densely dashed,rounded corners=3pt,inner sep=2mm},
  mainarrow/.style={->,thick,blue!65!black},
  monitor/.style={->,thick,blue!65!black,rounded corners=2pt},
  feedback/.style={->,thick,blue!65!black,rounded corners=3pt},
  influence/.style={->,densely dashed,gray!70!black},
  linklabel/.style={font=\footnotesize,align=center,fill=white,inner sep=1pt,text=blue!65!black}
]
  \tikzset{first/.style={},second/.style={},third/.style={}}
  \ifnum\policyfocus=1
    \tikzset{first/.style={added}}
  \fi
  \ifnum\policyfocus=2
    \tikzset{second/.style={added}}
  \fi
  \ifnum\policyfocus=3
    \tikzset{third/.style={added}}
  \fi

  \ifnum\policystage=1
    % 误差融合图按其自身内容收紧，不沿用风险分支所需的留白。
    \node[data,minimum width=34mm,text width=32mm] (img) at (-52mm,10mm) {多视角图像};
    \node[data,minimum width=34mm,text width=32mm] (ins) at (-52mm,-4mm) {子任务指令};
    \node[data,minimum width=34mm,text width=32mm,minimum height=17mm,first] (state) at (-52mm,-23mm)
      {增广状态\\本体状态与归一化\\视觉停靠误差};
    \node[process,minimum width=30mm,text width=28mm,minimum height=14mm] (vlm) at (-7mm,-4mm)
      {VLM主干网络\\编码观测与指令};
    \node[process,minimum width=30mm,text width=28mm,minimum height=13mm] (expert) at (-7mm,20mm)
      {动作专家\\条件流匹配};
    \node[container,fit=(expert)(vlm)] (pi) {};
    \node[linklabel,text width=16mm] at (-28mm,32mm) {策略模型\\\(\pi_{0.5}\)};
    \node[process,minimum width=30mm,text width=28mm,minimum height=17mm] (exec) at (44mm,-8mm)
      {机器人执行\\固定执行步数\(s\)\\再获取新观测};
    \node[offline,minimum width=58mm,text width=56mm,minimum height=18mm,first] (demo) at (-7mm,48mm)
      {随机停靠示范集\\固定位姿腕部观测，估计停靠误差\\融合误差输入，微调策略模型};

    \draw[mainarrow] (img.east) -- ([yshift=4mm]vlm.west);
    \draw[mainarrow] (ins.east) -- (vlm.west);
    \draw[mainarrow] (state.east) -- ([yshift=-4mm]vlm.west);
    \draw[mainarrow] (vlm.north) -- node[linklabel,right]{条件信息} (expert.south);
    \draw[mainarrow,rounded corners=2pt] (expert.east) --
      node[linklabel,above=1mm]{候选动作序列} (44mm,20mm) -- (exec.north);
    \draw[influence] (demo.south) -- (pi.north);
    \draw[feedback] (exec.south) -- (44mm,-40mm) --
      node[linklabel,below=1mm]{执行后的新观测} (-75mm,-40mm) -- (-75mm,10mm) -- (img.west);
  \else
  % 输入、VLM与动作专家保持原图的相对方位。
  \node[data] (img) at (-56mm,0mm) {多视角图像};
  \node[data] (ins) at (-56mm,-12mm) {子任务指令};
  \node[data,minimum width=34mm,text width=32mm,minimum height=18mm,first] (state) at (-55mm,-29mm)
    {增广状态\\本体状态与归一化\\视觉停靠误差};
  \node[process,minimum width=30mm,text width=28mm,minimum height=14mm] (vlm) at (-16mm,-12mm)
    {VLM主干网络\\编码观测与指令};
  \node[process,minimum width=30mm,text width=28mm,minimum height=14mm] (expert) at (-16mm,10mm)
    {动作专家\\条件流匹配};
  \node[container,fit=(expert)(vlm)] (pi) {};
  \node[linklabel,text width=16mm] at (-31mm,25mm) {策略模型\\\(\pi_{0.5}\)};

  \ifnum\policystage<3
    \node[process,minimum width=28mm,text width=26mm,minimum height=16mm] (exec) at (56mm,10mm)
      {机器人执行\\固定执行步数\(s\)\\再获取新观测};
  \else
    \node[process,minimum width=28mm,text width=26mm,minimum height=16mm] (exec) at (56mm,10mm)
      {机器人执行\\执行\(s_i\)步\\再获取新观测};
  \fi
  \node[offline,minimum width=64mm,text width=62mm,minimum height=17mm,first] (demo) at (-16mm,39mm)
    {随机停靠示范集\\固定位姿腕部观测，估计停靠误差\\融合误差输入，微调策略模型};

  \draw[mainarrow] (img.east) -- ([yshift=5mm]vlm.west);
  \draw[mainarrow] (ins.east) -- (vlm.west);
  \draw[mainarrow] (state.east) -- ([yshift=-5mm]vlm.west);
  \draw[mainarrow] (vlm.north) -- node[linklabel,right]{条件信息} (expert.south);
  \draw[influence] (demo.south) -- (pi.north);
  \draw[feedback] (exec.north) -- (56mm,54mm) --
    node[linklabel,above]{执行后的新观测} (-56mm,54mm) -- (img.north);

  % 基础动作通路先成立；调度方法在动作生成与下发之间引入RTC。
  \ifnum\policystage<3
    \draw[mainarrow] (expert.east) -- (exec.west);
    \node[linklabel] at (21mm,24mm) {候选动作序列};
  \else
    \node[process,minimum width=22mm,text width=20mm,minimum height=16mm,third] (rtc) at (23mm,10mm)
      {RTC引导\\生成时约束\\相邻序列};
    \draw[mainarrow] (expert.east) -- (rtc.west);
    \draw[mainarrow] (rtc.east) -- (exec.west);
    \node[linklabel,text width=20mm] at (5mm,25mm) {候选动作\\序列};
    \node[linklabel,text width=20mm] at (39mm,25mm) {引导后\\动作序列};
  \fi

  \ifnum\policystage>1
    % 保存上一轮序列；与当前未施加RTC约束的候选序列比较。
    \node[retained,minimum width=24mm,text width=22mm,minimum height=14mm,second] (tail) at (23mm,-12mm)
      {前一序列\\未执行尾部};
    \node[retained,minimum width=28mm,text width=26mm,minimum height=14mm,second] (difference) at (23mm,-37mm)
      {相邻序列\\预测偏差\(r_i\)};
    \coordinate (rawtap) at (5mm,10mm);
    \coordinate (historytap) at (39mm,10mm);
    \fill[blue!65!black] (rawtap) circle (0.65pt);
    \fill[blue!65!black] (historytap) circle (0.65pt);
    \draw[monitor] (historytap) -- (39mm,-12mm) -- (tail.east);
    \node[linklabel,text width=9mm,anchor=west] at (40mm,-4mm) {上轮\\尾部};
    \draw[monitor] (tail.south) -- node[linklabel,right]{对齐比较} (difference.north);
    \draw[monitor] (rawtap) -- (5mm,-37mm) -- (difference.west);
    \fill[blue!65!black] (5mm,-37mm) circle (0.65pt);

    % 五项输入均保留可追溯的入边，预测偏差只是风险输入之一。
    \node[retained,minimum width=62mm,text width=60mm,minimum height=22mm,second] (risk) at (-7mm,-60mm)
      {轻量风险判别模型\\
       \(f_{\boldsymbol\psi}(\mathbf z_i,\widehat{\mathbf A}^{\mathrm{raw}}_i,
       \tilde{\mathbf x}_{t_i},\ell_{t_i},r_i)\)\\
       输出：失败风险分数\(\rho_i\)};
    \node[retained,minimum width=28mm,text width=26mm,minimum height=22mm,second] (dec) at (56mm,-60mm)
      {提前停止判定\\研究预警及时性\\与误报的权衡\\满足条件则停止};
    \node[offline,minimum width=62mm,text width=60mm,minimum height=14mm,second] (rtrain) at (-7mm,-83mm)
      {误差注入回合与失败标签\\冻结\(\pi_{0.5}\)，仅训练风险模型};

    \draw[monitor] (vlm.south) -- node[linklabel,left,pos=0.55]{观测特征\(\mathbf z_i\)}
      (vlm.south |- risk.north);
    \draw[monitor] (state.south) -- (-55mm,-55mm) --
      node[linklabel,above]{增广状态} ([yshift=5mm]risk.west);
    \draw[monitor] (ins.west) -- (-74mm,-12mm) -- (-74mm,-67mm) --
      node[linklabel,above]{子任务指令} ([yshift=-7mm]risk.west);
    \draw[monitor] (5mm,-37mm) --
      node[linklabel,left,pos=0.62]{候选动作\\序列}
      (5mm,-37mm |- risk.north);
    \draw[monitor] (difference.south) -- ([xshift=25mm]risk.north);
    \draw[monitor] (risk.east) -- node[linklabel,above]{风险分数} (dec.west);
    \draw[influence] (rtrain.north) -- (risk.south);

    \ifnum\policystage=2
      \draw[monitor] (dec.north) -- node[linklabel,text width=20mm,pos=0.48]{停止或放行} (exec.south);
    \else
      % 风险分数调度实际步数，停止路径始终单独保留。
      \node[retained,minimum width=28mm,text width=26mm,minimum height=19mm,third] (schedule) at (56mm,-29mm)
        {风险驱动调度\\低风险长步数\\高风险短步数};
      \draw[monitor] (dec.north) -- node[linklabel,text width=20mm]{放行及风险} (schedule.south);
      \draw[monitor] (schedule.north) -- node[linklabel,text width=18mm,pos=0.23]{执行步数\(s_i\)} (exec.south);
      \draw[monitor] (dec.east) -- (73mm,-60mm) --
        node[linklabel,rotate=90]{提前停止} (73mm,10mm) -- (exec.east);
      \draw[monitor] (tail.north) -- (rtc.south);
    \fi
  \fi
  \fi
\end{tikzpicture}
